\section*{Capítulo \ref{ch:b1:chemical-biogenic-sediments} --- Sedimentos carbonáticos, silicosos e evaporíticos}

\begin{solution}{exo:b1:chemical-biogenic-sediments:1}
Cloreto $0.5503\times35.165 = \qty{19.35}{g}$; sódio $0.3066\times35.165 =
\qty{10.78}{g}$; sulfato \qty{2.711}{g}; magnésio \qty{1.284}{g}. Sódio e
cloreto juntos somam $0.5503 + 0.3066 = 0.857$, isto é \qty{85.7}{\%} do
sal.
\end{solution}

\begin{solution}{exo:b1:chemical-biogenic-sediments:2}
A fração que resta é $1/\mathrm{CF}$: carbonato de cálcio $1/1.8 = 0.556$ restante,
\qty{44.4}{\%} evaporado; gipsita $0.263$ restante, \qty{73.7}{\%} evaporado;
halita $0.094$ restante, \qty{90.6}{\%} evaporado; sulfato de magnésio $0.014$ restante,
\qty{98.6}{\%} evaporado.
\end{solution}

\begin{solution}{exo:b1:chemical-biogenic-sediments:3}
(a)~$\rho_s(1-\phi) = G/a = 8.4/0.0042 = \qty{2000}{kg.m^{-3}}$. (b)~Com
aragonita, $1 - 2000/2930 = 0.32$, uma porosidade de \qty{32}{\%}; com calcita,
$1 - 2000/2710 = 0.26$. (c)~$a = 3.7/2000 = \qty{1.85e-3}{m}$, \qty{1.85}{mm} por
ano.
\end{solution}

\begin{solution}{exo:b1:chemical-biogenic-sediments:4}
Até \qty{4000}{m} o teor cai apenas devagar (88 a 80); abaixo de
\qty{4000}{m} ele cai bruscamente: a \omterm{def:b1:chemical-biogenic-sediments:ccd}{lisoclina} fica a cerca de \qty{4000}{m}. O
teor cai através de \qty{10}{\%} entre \qty{4700}{m} (12) e \qty{4800}{m}
(6): $4700 + 100\times(12-10)/(12-6) = \qty{4733}{m}$, uma CCD a cerca de
\qty{4730}{m}, cerca de \qty{700}{m} abaixo da \omterm{def:b1:chemical-biogenic-sediments:ccd}{lisoclina}.
\end{solution}

\begin{solution}{exo:b1:chemical-biogenic-sediments:5}
(a)~\ce{Ca^2+ + 2 HCO3- -> CaCO3 + CO2 + H2O}: um mol de \ce{CO2} por mol de
\ce{CaCO3}, portanto $44.01/100.09 = \qty{0.440}{kg}$ de \ce{CO2} por quilograma de
\ce{CaCO3}. (b)~$3.7\times0.440 = \qty{1.63}{kg}$ de \ce{CO2} por metro quadrado e
por ano. (c)~Dos dois íons hidrogenocarbonato usados, apenas um carbono entra no
\omterm{def:b1:rock-forming-minerals:mineral}{mineral}; o outro sai como gás carbônico. A calcificação retira alcalinidade da água e
reduz sua capacidade de reter carbono dissolvido; quanto do \ce{CO2} chega ao ar
depende do tamponamento da água do mar, uma questão dos equilíbrios carbonáticos dos
livros de química.
\end{solution}

\begin{solution}{exo:b1:chemical-biogenic-sediments:6}
(a)~$t_c = (2000/350)^2 = \qty{32.7}{Myr}$; a \qty{25}{km} por milhão de anos,
$25\times32.7 = \qty{816}{km}$ do eixo. (b)~$t_c = (2600/350)^2 =
\qty{55.2}{Myr}$, \qty{1380}{km}. (c)~Um fundo de \qty{60}{Myr} recebeu argila
vermelha por $60 - 32.7 = \qty{27.3}{Myr}$ no primeiro caso, e por apenas
$60 - 55.2 = \qty{4.8}{Myr}$ no segundo.
\end{solution}

\begin{solution}{exo:b1:chemical-biogenic-sediments:7}
(a)~Sim: $a = 4.2 > r = 3$, e o \omterm{def:b1:chemical-biogenic-sediments:reef}{recife} ganha \qty{1.2}{mm} por ano sobre o mar
até que seu topo chegue à superfície. (b)~$D$ aumenta $10 - 4.2 = \qty{5.8}{mm}$ por
ano; de 2 a \qty{50}{m} leva $48/0.0058 = \num{8300}$ anos. (c)~Porque
sua taxa de crescimento cai à medida que a água se aprofunda: chega menos luz às algas
que vivem nos corais e os alimentam, de modo que $a$ diminui e o \omterm{def:b1:chemical-biogenic-sediments:reef}{recife} fica para trás
cada vez mais depressa.
\end{solution}

\begin{solution}{exo:b1:chemical-biogenic-sediments:8}
(a)~A reação é \ce{CaSO4.2H2O -> CaSO4 + 2 H2O}. (b)~Os volumes molares são $172.17/2.30 =
\qty{74.9}{cm^3}$ e $136.14/2.97 = \qty{45.8}{cm^3}$, de modo que o sólido perde
$1 - 45.8/74.9 = \qty{38.8}{\%}$ de seu volume. (c)~$2\times18.02/172.17 = 0.209$:
\qty{20.9}{\%} da massa sai como água. Liberada em profundidade, essa água eleva a pressão
nos poros das rochas ao redor; sob um selo ela pode enfraquecer as rochas e ajudar falhas
a deslizar, e pode dissolver e transportar outros sais.
\end{solution}

\begin{solution}{exo:b1:chemical-biogenic-sediments:9}
(a)~$9.2/255 = \qty{3.6}{\%}$ preservado, \qty{96.4}{\%} dissolvido. (b)~Onde a
produção é muito alta, algumas por cento dela ainda são um fluxo grande, e onde pouca
argila ou carbonato a dilui: sob as águas de ressurgência do oceano Antártico, do Pacífico
norte e do Pacífico equatorial, e abaixo da CCD. (c)~Carapaças finas e delicadas se
dissolvem primeiro e as robustas sobrevivem: o sedimento guarda uma seleção das espécies
mais resistentes, e não a comunidade que vivia na água.
\end{solution}

\begin{solution}{exo:b1:chemical-biogenic-sediments:10}
(a)~$EA = 1.5\times\qty{1e10}{m^2} = \qty{1.5e10}{m^3}$ por ano. (b)~Com
$Q_{\mathrm{in}} = \qty{1.8e10}{m^3}$, $Q_{\mathrm{out}} = \qty{0.3e10}{m^3}$ e
$k = 6$: acima do limiar da gipsita (3.8), mas não do da halita (10.6), então
carbonato de cálcio e gipsita precipitam. Com $Q_{\mathrm{in}} =
\qty{1.575e10}{m^3}$, $Q_{\mathrm{out}} = \qty{0.075e10}{m^3}$ e $k = 21$, para
um íon que permanece em solução, como o magnésio: a salmoura passa do limiar da halita,
então a halita também precipita, mas não os sais de magnésio (70). (c)~A halita
exige $Q_{\mathrm{in}}/(Q_{\mathrm{in}} - EA) \ge 10.6$, isto é,
$Q_{\mathrm{in}} \le 10.6\,EA/9.6 = 1.104\,EA = \qty{1.66e10}{m^3}$ por ano: a
entrada pode exceder a evaporação em no máximo cerca de um décimo.
\end{solution}

\begin{solution}{exo:b1:chemical-biogenic-sediments:11}
(a)~A água do mar contém \qty{67.3}{mg} de brometo por quilograma; a 10.6~vezes, a
salmoura contém $67.3\times10.6 = \qty{713}{mg.kg^{-1}}$ e a primeira halita
$0.032\times713 = \qty{22.8}{mg.kg^{-1}}$. (b)~A 30~vezes, $67.3\times30 =
\qty{2019}{mg.kg^{-1}}$ na salmoura e \qty{64.6}{mg.kg^{-1}} na halita.
(c)~A halita retira apenas 3.2~centésimos da concentração da salmoura, de modo que
quase todo o brometo permanece em solução e continua se concentrando à medida que a
água evapora. (d)~A subida registra a concentração da salmoura; a queda brusca registra
sua diluição, por uma nova entrada de água do mar na bacia, ou de água doce, antes que
a halita continuasse a crescer.
\end{solution}

\begin{solution}{exo:b1:chemical-biogenic-sediments:12}
(a)~$t_c = \qty{32.65}{Myr}$, então $0.02\times\num{32650} = \qty{653}{m}$ de
lodo (\num{32650} intervalos de mil anos a \qty{2}{cm}) e
$0.001\times\num{37350} = \qty{37}{m}$ de argila vermelha. (b)~$t_c = (2500/350)^2 =
\qty{51.0}{Myr}$: \qty{1020}{m} de lodo e \qty{19}{m} de argila vermelha. Uma CCD
mais funda dá um lodo muito mais espesso, e a argila vermelha do fundo do Atlântico
profundo é mais fina que a do Pacífico da mesma idade. (c)~A CCD pode ter sido mais funda
por algum tempo, de modo que o local voltou a subir acima dela; ou a placa levou o local
sob um cinturão de alta produtividade, como o equador, onde mais carbonato chovia; ou a
camada foi trazida por um deslizamento ou uma \omterm{def:b1:clastic-sediments:turbidite}{corrente de turbidez} de uma parte mais alta
do fundo do mar.
\end{solution}

\begin{solution}{pb:b1:chemical-biogenic-sediments:1}
\textbf{1.} $1030\times1000 = \qty{1.03e6}{kg}$ de água do mar, contendo
$1.03\times10^6\times0.035165 = \qty{36200}{kg}$ de sal.

\textbf{2.} Gipsita $1.03\times10^6\times1.581\times10^{-3} = \qty{1628}{kg}$;
halita $1.03\times10^6\times0.02741 = \qty{28230}{kg}$.

\textbf{3.} Carbonato de cálcio $113.3/2710 = \qty{0.042}{m}$; gipsita $1628/2300 =
\qty{0.71}{m}$; halita $28230/2170 = \qty{13.0}{m}$.

\textbf{4.} A anidrita tem massa $1628\times136.14/172.17 =
\qty{1288}{kg}$ e espessura $1288/2970 = \qty{0.43}{m}$.

\textbf{5.} Seco, o sal é $0.110 + 1.250 + 27.41 + 6.34 = \qty{35.11}{g}$ por
quilograma, dos quais a halita é $27.41/35.11 = \qty{78}{\%}$: sódio e cloreto
formam a maior parte do sal dissolvido, e a halita retira quase todos eles.

\textbf{6.} De baixo para cima: carbonato de cálcio, gipsita, halita, os sais amargos.
Cada \omterm{def:b1:rock-forming-minerals:mineral}{mineral} começa a precipitar quando a salmoura atinge sua saturação, o menos solúvel
primeiro, e se deposita sobre o que foi depositado antes.

\textbf{7.} A profundidade da salmoura é $1000/\mathrm{CF}$: \qty{556}{m} para carbonato de
cálcio, \qty{263}{m} para gipsita, \qty{94}{m} para halita.

\textbf{8.} $1 - 1/10.6 = \qty{90.6}{\%}$.

\textbf{9.} $1000/70 = \qty{14}{m}$ para sulfato de magnésio, $1000/90 = \qty{11}{m}$
para os sais de potássio.

\textbf{10.} Cerca de \qty{1.10}{g.cm^{-3}} quando se forma carbonato de cálcio, 1.13 com
gipsita, 1.22 com halita e mais de 1.30 com os sais amargos, contra 1.03 para a água do
mar. A água do mar que entra flutua sobre a salmoura densa, que permanece no fundo; em uma
bacia barrada a salmoura pode retornar ao oceano abaixo da entrada.

\textbf{11.} É apenas de \qty{4}{cm} de espessura, contra \qty{13}{m} de halita; e forma-se
nas margens, onde a água que entra se concentra primeiro, mais do que no centro, de modo que
é misturada à base da gipsita ou substituída por ela.

\textbf{12.} A halita retira apenas uma pequena parte do brometo da salmoura
(exercício~\ref{exo:b1:chemical-biogenic-sediments:11}), de modo que a salmoura fica mais
rica em brometo à medida que a halita cresce, e cada nova camada de halita incorpora mais.

\textbf{13.} $1500/13.01 = 115.3$: 116 enchimentos.

\textbf{14.} $116\times0.434 = \qty{50}{m}$ de anidrita, não como um só leito, mas como
116 camadas, cada uma um pouco menos de meio metro de espessura, na base da halita de
seu enchimento: uma sucessão cíclica de anidrita e halita.

\textbf{15.} Por enchimento, $1.03\times10^6\times6.34\times10^{-3}/2000 =
\qty{3.3}{m}$, de modo que \qty{379}{m} para 116 enchimentos. São os sais mais solúveis: a
entrada seguinte de água do mar os redissolve, e uma bacia cuja salmoura sai antes do fim
nunca os alcança.

\textbf{16.} A cada reenchimento a salmoura é diluída: o teor de brometo cai bruscamente na
base da nova halita e sobe de novo para cima, um perfil em dente de serra com um dente por
enchimento.

\textbf{17.} $2000/13.01 = 154$ enchimentos.

\textbf{18.} $5000\times365.25\times86400 = \qty{1.578e11}{m^3}$ por ano, distribuídos sobre
$\qty{1.0e11}{m^2}$: $E = \qty{1.58}{m}$ por ano.

\textbf{19.} $1.578\times10^{11}\times1030\times0.035165 = \qty{5.72e12}{kg}$ de sal por ano.

\textbf{20.} $1.578\times10^{11}\times1030\times0.02741/2170 =
\qty{2.05e9}{m^3}$ de halita por ano, isto é, \qty{0.0205}{m}, cerca de \qty{2}{cm} por ano
sobre o fundo.

\textbf{21.} $1500/0.0205 = \num{73000}$ anos.

\textbf{22.} O nível cai \qty{1.58}{m} por ano; \qty{1000}{m} de água se vão em
$1000/1.578 = 634$ anos.

\textbf{23.} A halita aparece quando \qty{90.6}{\%} da água já se foram, após
$905.7/1.578 = 574$ anos.

\textbf{24.} Os \qty{1.5}{km} de halita levaram o sal de 116 enchimentos da bacia
(115.3, a arredondar para cima), quer tenham vindo um após o outro à medida que a bacia
secava e reenchia, quer continuamente por cima da soleira, como na questão~21, à taxa de
uma bacia de água do mar a cada $1000/1.578 = 634$ anos.
\end{solution}
